\documentclass{article}
\usepackage[T1]{fontenc}
\usepackage{lmodern}
\usepackage{graphicx}
\usepackage{amsmath,amssymb}
\usepackage[a4paper,margin=2cm]{geometry}
\usepackage[absolute,overlay]{textpos}
\setlength{\TPHorizModule}{1mm}
\setlength{\TPVertModule}{1mm}
\usepackage[indonesian]{babel}
\usepackage{hyperref}
\setlength{\emergencystretch}{2em}
% unit: Halaman 11

\begin{document}

% === Halaman 11 (llm) ===

\section*{Sejarah Steganografi}

\begin{itemize}
    \item Usia steganografi setua usia kriptografi, dan sejarah keduanya berjalan bersamaan.
\end{itemize}

\begin{itemize}
    \item Periode sejarah steganografi dapat dibagi menjadi:
    \begin{enumerate}
        \item Steganografi kuno (\textit{ancient steganography})
        \item Steganografi zaman renaisans (\textit{renaissance steganography})
        \item Steganografi zaman perang dunia
        \item Steganografi modern
    \end{enumerate}
\end{itemize}

\end{document}
